\documentclass[11pt]{article}
\usepackage{lecturenotes}

% Uncomment to change the colour of comments in code for this document only.
% \definecolor{codecomment}{gray}{0.25}

\course{CPSC 221: Basic Algorithms and Data Structures}

\begin{document}

\lecture{N}{Lecture Title}{Month Day, 2026}

\section{Introduction}

Inline code looks like \cppinline{std::vector<int>}, and a short snippet
looks like this:

\begin{cpp}
#include <vector>

// Returns the sum of the entries of v.
int sum(const std::vector<int>& v) {
    int total = 0;
    for (int x : v) {
        total += x; /* accumulate */
    }
    return total;
}
\end{cpp}

\section{Algorithms}

\begin{algorithm}[label=alg:findmin]{Finding the minimum of an array}
// Precondition: n >= 1.
int findMin(const int a[], int n) {
    int best = a[0];
    for (int i = 1; i < n; i++) {    (*@\label{line:loop}@*)
        if (a[i] < best) {
            best = a[i];             // new smallest so far
        }
    }
    return best;
}
\end{algorithm}

\Cref{alg:findmin} makes $n - 1$ comparisons, one per iteration of the loop
on line~\ref{line:loop}.

\section{Theorem environments}

\begin{definition}[Big-O]\label{def:bigo}
  $f(n) \in O(g(n))$ if there exist constants $c > 0$ and $n_0$ such that
  $f(n) \le c \cdot g(n)$ for all $n \ge n_0$.
\end{definition}

\begin{theorem}[Correctness of \texttt{findMin}]\label{thm:findmin}
  For every array \cppinline{a} of length $n \ge 1$, \cref{alg:findmin}
  returns $\min_{0 \le i < n} a[i]$.
\end{theorem}

\begin{proof}
  By induction on the loop invariant: at the start of iteration $i$,
  \cppinline{best} $= \min_{0 \le j < i} a[j]$.
\end{proof}

\begin{lemma}\label{lem:example}
  A lemma.
\end{lemma}

\begin{proposition}
  A proposition.
\end{proposition}

\begin{corollary}
  A corollary of \cref{thm:findmin} and \cref{lem:example}.
\end{corollary}

\begin{claim}
  A claim.
\end{claim}

\begin{fact}
  A fact.
\end{fact}

\begin{conjecture}
  A conjecture.
\end{conjecture}

\begin{example}
  An example using \cref{def:bigo}: $3n + 5 \in O(n)$.
\end{example}

\begin{exercise}
  An exercise.
\end{exercise}

\begin{problem}
  A problem.
\end{problem}

\begin{theorem*}
  An unnumbered theorem (every environment has a starred version).
\end{theorem*}

\begin{remark}
  Remarks and notes are not boxed.
\end{remark}

\begin{note}
  A note.
\end{note}

\end{document}
